% Viva question bank, \input by SonicVitals_Viva_Guide.tex
\newpage
\section{Viva question bank}
{\small Short answers you can say out loud. \textbf{Bold} = the key point. Where a question is a trap, the answer says so.}

\newcommand{\qa}[2]{\par{\raggedright\bfseries Q.\ #1\par}\nopagebreak #2\par\vspace{3pt}}
\newcommand{\qgroup}[1]{\vspace{2pt}{\normalsize\bfseries\color{acc}#1}\par\vspace{1pt}}

\begin{multicols}{2}
\footnotesize
\setlength{\columnsep}{14pt}

% ---------------------------------------------------------------
\qgroup{A. Motivation and design choices}

\qa{Why not just use a phone or a webcam? The ESP32 image is poor.}
{\textbf{Not for image quality} (a phone is better): for (1)~\textbf{one clock} for camera and mic, so fusion is time-aligned;
(2)~\textbf{sensor control}: we lock exposure, gain and white balance after 3\,s, which phones and browser webcams keep
changing; (3)~a \textbf{chest microphone} for heart sounds, which a phone filming your face can't provide; (4)~a
\textbf{cheap, standalone node} with its own display and Wi-Fi. Trade-off accepted: 240$\times$160, about 8\,fps. The same software
also runs on a webcam (fallback mode), so a direct comparison is possible.}

\qa{What problem does the project solve?}
{Measuring vital signs \textbf{without touching} the person (no finger clip or chest strap): useful for burns, neonates, infection
control and continuous monitoring. We add a second modality (heart sounds) to check the camera estimate.}

\qa{What is novel here?}
{Not the algorithms (CHROM, Welch, PCG envelopes are published). The contribution is the \textbf{integration}: a
low-cost ESP32 node that streams time-synchronised video and audio, a camera/acoustic \textbf{fusion with an agreement gate},
\textbf{quality gating} that shows ``--'' instead of noise, and an honest study of what such a cheap sensor can and cannot do.}

\qa{Why ``SonicVitals''?}
{``Sonic'' for the acoustic (heart-sound) channel and ``Vitals'' for HR/RR; the system is multimodal: optical plus acoustic.}

\qa{Who would use it?}
{A research or screening prototype: telehealth kiosks, hostels and labs, low-resource clinics. \textbf{Not a medical device}; it
would need clinical validation and regulatory approval.}


% ---------------------------------------------------------------
\qgroup{B. Hardware}

\qa{Why the ESP32-S3?}
{Dual-core 240\,MHz Xtensa LX7, \textbf{8\,MB PSRAM} (frame buffer), a camera (DVP) interface, an \textbf{I$^2$S} interface for the
mic, built-in \textbf{Wi-Fi}, low cost. Two cores let audio and camera run in parallel without blocking.}

\qa{Which camera, and what is its limitation?}
{\textbf{GC2145} (2\,MP CMOS). It has \textbf{no hardware JPEG encoder}: the driver rejected JPEG mode, so we capture RGB565 and
compress in software (quality 60, about 5\,kB per frame). That limits us to about 8\,fps over Wi-Fi.}

\qa{Why the INMP441 microphone?}
{A digital I$^2$S MEMS mic: 24-bit output, so no analogue amplifier or ADC noise, and it is cheap and small.
\textbf{Limitation:} its response falls off below about 60\,Hz, while heart sounds are 20--150\,Hz, so it must be pressed on the
chest; in air the heart sounds are weak.}

\qa{What is the OLED for?}
{An SSD1306 128$\times$64 display over I$^2$C (GPIO47 SDA / 48 SCL, address 0x3C). It shows HR/RR/status sent back by the host and
the node's IP address, so the device is usable without the dashboard.}

\qa{Why couldn't you use USB for everything?}
{The board has one USB port, wired to a CH340 USB-to-serial chip on \textbf{GPIO43/44}. The mic's I$^2$S clock lines were on
the same pins, so with the mic on, the serial link died (verified). Also the CH340 was unreliable above 1\,Mbaud. So we stream
over \textbf{Wi-Fi}; USB is only for flashing and power.}

\qa{How is it powered? Power consumption?}
{5\,V over USB (laptop or power bank). \textbf{Power was not yet measured}; planned with a USB power meter (Table III of the paper). Do not
quote a number you haven't measured.}

\qa{Why PSRAM?}
{A frame ring buffer (24 slots $\times$ 16\,kB $\approx$ 3\,s) absorbs Wi-Fi stalls, so frames are not lost when the hotspot
briefly slows down.}


% ---------------------------------------------------------------
\qgroup{C. Firmware and communication}

\qa{How is the firmware organised?}
{FreeRTOS tasks: \textbf{audio task on core 0} (reads I$^2$S DMA, filters, queues blocks), \textbf{camera task on core 1}
(captures at a fixed 8\,fps, JPEG-encodes, puts frames in the PSRAM ring), and the main loop sends data, reads host commands and
updates the OLED.}

\qa{How are camera and audio synchronised?}
{Every frame and audio block carries a \textbf{microsecond timestamp from the same ESP32 clock} (\texttt{esp\_timer}) plus a
sequence number. The host analyses by timestamp, not arrival time, so network delay doesn't matter.}

\qa{Why TCP and not UDP?}
{A lost UDP packet would corrupt a JPEG frame. TCP guarantees order and delivery. Latency isn't critical because data is
timestamped at capture.}

\qa{Why not Bluetooth?}
{BLE throughput (tens of kB/s in practice) is too low for video plus audio; Wi-Fi gives about 70\,kB/s on a phone hotspot,
which is enough.}

\qa{What does the protocol look like?}
{Text lines: \texttt{<<HELLO:...>>}, \texttt{<<FRAME\_START:seq,t\_us,len>>} + base64 JPEG + \texttt{<<FRAME\_END>>},
\texttt{<<AUDIO:seq,t\_us,n,data>>}, \texttt{<<STAT:...>>} (fps, drops). Downlink \texttt{<<VITALS:HR,RR,SpO2,Q,STATUS>>} to the
OLED. The length field and sequence numbers detect corrupt or missing data.}

\qa{Why base64, which adds 33 \%?}
{It keeps the stream as plain text lines, so the parser can always re-synchronise at the next \texttt{<<} marker and it is easy to
debug. The cost is $4/3$ bandwidth; binary framing is future work.}

\qa{Why send JPEG and not raw pixels?}
{Raw: $240\times160\times2 = 76.8$\,kB per frame, 614\,kB/s at 8\,fps, far above the link. JPEG is about 5\,kB.
Downside: compression removes some of the tiny colour changes (McDuff et al., 2017).}

\qa{Why decimate audio from 16 kHz to 4 kHz?}
{Bandwidth: 16\,kHz $\times$ 2 bytes $=32$\,kB/s (43\,kB/s in base64) would starve the video. Heart sounds are below 400\,Hz, so
4\,kHz (Nyquist 2\,kHz) keeps them. A \textbf{63-tap FIR low-pass} filter runs first to prevent \textbf{aliasing}.}

\qa{Why lock exposure, gain and white balance?}
{Auto-exposure changes brightness every few frames, a change much larger than the pulse (0.1--1\,\%). After 3\,s of settling,
we freeze them so colour changes come from the skin, not the camera.}

\qa{How does the laptop find the board?}
{The node joins the phone hotspot and announces \texttt{sonicvitals.local} (mDNS); the host also scans the /24 subnet for port
3333. If no network is configured, the node becomes an access point (\texttt{SonicVitals-Node}, 192.168.4.1).}

\qa{How did you verify the link quality?}
{Measured over 30\,s: \textbf{8.0\,fps}, frame-timing jitter \textbf{0.4\,ms}, \textbf{0} dropped or corrupt frames, \textbf{0} audio
blocks lost (sequence numbers prove it).}


% ---------------------------------------------------------------
\qgroup{D. rPPG theory and signal processing}

\qa{What is photoplethysmography?}
{Measuring blood-volume changes optically. In a pulse oximeter a LED shines \emph{into} the finger; \textbf{remote} PPG uses
ambient light and a camera at a distance.}

\qa{Why does skin colour change with the heartbeat?}
{Each beat pushes blood into the capillaries; haemoglobin \textbf{absorbs green light strongly} (less red), so reflected green
dips slightly in systole. The change is about 0.1--1\,\% of the brightness.}

\qa{How can a change below 1 grey level be measured?}
{By \textbf{averaging thousands of pixels}: random noise falls as $1/\sqrt{N}$, while the pulse is common to all skin pixels and
survives.}

\qa{Why the forehead and cheeks?}
{Thin skin, dense capillaries, little hair, and relatively still. Forehead plus both cheeks gives more pixels than the forehead
alone, so lower noise. Beard and eyebrows are avoided.}

\qa{Explain CHROM.}
{Normalise R, G, B by their means (removes light intensity). Build $X=3R_n-2G_n$ and $Y=1.5R_n+G_n-1.5B_n$; both contain the
motion/specular part, so $S=X-\alpha Y$ with $\alpha=\sigma_X/\sigma_Y$ \textbf{cancels it} and keeps the pulse (de Haan \& Jeanne
2013).}

\qa{Why not just the green channel?}
{Green has the strongest pulse but also all the lighting and motion changes. CHROM is more robust; we keep green as a baseline
for comparison (ablation).}

\qa{CHROM vs POS vs deep learning?}
{POS (Wang 2017) is a similar projection; either works. Deep models (DeepPhys, PhysNet) can be more accurate but need large
training sets and a GPU, and may not generalise to a new low-quality camera. CHROM needs no training.}

\qa{Why band-pass 0.75--2.5 Hz?}
{0.75--2.5\,Hz $=$ 45--150\,BPM, the expected range. It removes breathing ($<0.4$\,Hz) and high-frequency noise.
\textbf{Limitation}: HR below 45 (athletes) would be missed.}

\qa{Why Butterworth, and why filtfilt?}
{Butterworth has a \textbf{maximally flat passband} (no ripple distorting amplitudes). \texttt{filtfilt} runs it forwards and
backwards, giving \textbf{zero phase shift}, so peaks are not delayed.}

\qa{Why resample to 15 Hz?}
{Frames arrive at uneven times over Wi-Fi; the FFT needs uniform samples. Each frame is placed at its capture timestamp and
linearly interpolated onto a uniform grid.}

\qa{Is 8 fps enough? (Nyquist)}
{The sampling rate must exceed twice the highest frequency. HR up to 150\,BPM is 2.5\,Hz, which needs $>5$\,fps; 8\,fps gives
4\,Hz Nyquist. Enough for HR, not for fine waveform shape or HRV.}

\qa{What is Welch's method and why use it?}
{Split the signal into 50\,\%-overlapping segments, apply a Hann window, average the periodograms. Averaging \textbf{reduces
the variance} (noise) of the spectrum compared with a single FFT.}

\qa{Why a Hann window?}
{Cutting a signal abruptly causes \textbf{spectral leakage} (power smeared into other bins). The Hann window tapers the ends and
reduces leakage.}

\qa{What is the frequency resolution?}
{$\Delta f=1/T$. A 10\,s window gives 0.1\,Hz $=6$\,BPM. \textbf{Zero-padding} (FFT length 4096) only interpolates the spectrum;
\textbf{parabolic interpolation} through the peak and its neighbours then estimates the peak between bins.}

\qa{Why wait 10 s before showing HR?}
{To get about 8--15 beats and a usable frequency resolution. Shorter windows give unstable peaks. The analysis buffer is 30\,s.}

\qa{How do you know a reading is good? (SNR)}
{The de Haan \textbf{SNR}: power near $f_0$ and its harmonic $2f_0$ divided by the rest of 0.5--4\,Hz, in dB. HR is accepted only if
SNR $\geq 0$\,dB. We set this threshold by testing pure noise: it passes in fewer than 5\,\% of cases.}

\qa{What is spectral concentration $Q$?}
{The share of in-band power within $\pm0.1$\,Hz of the peak. High $Q$ means one clear peak. It is used as the camera's fusion
weight.}

\qa{What happens when the person moves briefly?}
{Gaps up to \textbf{2\,s} are bridged by interpolation, so the estimate continues. Longer breaks restart the 10\,s buffer while the
last value is \textbf{held} (marked ``HELD'') for up to 15\,s.}

\qa{How do you stop the HR jumping between peaks?}
{\textbf{Peak tracking}: if a peak within $\pm12$\,BPM of the last accepted HR has at least half the power of the highest peak,
it is kept. The displayed value is a median of the last 3 estimates followed by an exponential moving average (EMA).}

\qa{How does face detection work?}
{\textbf{Viola--Jones Haar cascade}: Haar-like rectangle features computed quickly from an integral image, chosen by AdaBoost,
arranged in a cascade that rejects non-faces early. Fast on a CPU. The box is smoothed to reduce jitter.}

\qa{Why not a deep face detector (MediaPipe)?}
{Haar is light and real-time on a laptop CPU. MediaPipe Face Mesh would track landmarks better under head turns: \textbf{future
work}.}

\qa{What are the five face-gate tests?}
{Not touching the edges; size 20--88\,\% of the frame width; aspect ratio 0.65--1.30 (not turned); all ROIs inside the frame;
enough light and similar left/right brightness. Frames that fail are not used.}

\qa{Two faces in the frame?}
{The largest face is used. Multi-person tracking is future work.}


% ---------------------------------------------------------------
\qgroup{E. Respiration, SpO$_2$ and HRV}

\qa{How do you get the breathing rate from a face?}
{Breathing modulates the skin signal (venous return, chest and head movement, light angle): \textbf{respiratory-induced intensity
variation}. Band-pass green 0.15--0.40\,Hz (9--24 breaths/min), take the spectral peak. Needs 15\,s and a clear peak
($\geq 50\,\%$ of the low-frequency power), otherwise ``--''.}

\qa{Can a camera measure SpO$_2$?}
{Only roughly. Pulse oximeters use the \textbf{ratio of ratios} $(AC/DC)_{red}/(AC/DC)_{IR}$; our camera has no infrared, so we
use red/blue as a proxy. It needs \textbf{calibration} against a reference oximeter ($\text{SpO}_2=A-B\cdot R$). Without that file,
the dashboard shows only the uncalibrated ratio.}

\qa{What is HRV and how do you measure it?}
{Beat-to-beat variation of the inter-beat interval (IBI). RMSSD $=\sqrt{\text{mean}(\Delta IBI^2)}$. We compute it from heart-sound
beat times and need $\geq 20$ clean intervals. At 8\,fps the camera's time resolution (125\,ms) is too coarse for HRV.}

\qa{What is LF/HF?}
{The ratio of HRV power at 0.04--0.15\,Hz to 0.15--0.4\,Hz, used as a sympathovagal-balance (stress) index. It needs several
minutes of accurate beats, so we \textbf{removed} the stress card (it was a placeholder).}


% ---------------------------------------------------------------
\qgroup{F. Heart sounds (PCG)}

\qa{What are S1 and S2?}
{S1 (``lub''): mitral and tricuspid valves closing at the start of systole. S2 (``dub''): aortic and pulmonary valves closing at the
end. Main energy is about 20--150\,Hz.}

\qa{How do you get HR from sound?}
{Band-pass 25--400\,Hz, resample to 1\,kHz, \textbf{Shannon energy} envelope $-x^2\log x^2$ smoothed over 20\,ms, then
\textbf{autocorrelation}: the first strong peak at lag $\tau$ (0.4--1.33\,s) is one beat period, so $\text{HR}=60/\tau$.}

\qa{Why Shannon energy and not just $x^2$?}
{$x^2$ is dominated by the loudest peaks. Shannon energy emphasises \textbf{medium-intensity} sounds and suppresses both low
noise and very high spikes, giving a cleaner envelope.}

\qa{Why autocorrelation instead of peak picking?}
{It uses the \textbf{periodicity of the whole window}, so it is robust to a few missed or extra peaks; it does not need to label
S1 vs S2.}

\qa{What is $\rho$?}
{The autocorrelation value at the beat lag (0--1): how periodic the sound is. Below 0.3, the acoustic HR is not reported.}

\qa{Did the heart-sound channel work?}
{\textbf{Honestly: weakly.} In air, 99.5\,\% of the audio power was below 25\,Hz and $\rho$ stayed $\leq 0.18$, because the
INMP441 cuts off around 60\,Hz. It needs firm chest contact or a stethoscope-coupled mic: \textbf{future work}.}


% ---------------------------------------------------------------
\qgroup{G. Fusion}

\qa{How do you fuse camera and sound?}
{Quality-weighted mean $(w_oHR_o+w_aHR_a)/(w_o+w_a)$, with $w_o=Q$ (camera) and $w_a=\rho$ (sound). \textbf{Agreement gate:} only if
they differ by $\leq 8$\,BPM; otherwise use the better channel and flag ``CONFLICT''.}

\qa{Why an agreement gate?}
{Averaging a correct value with a wrong one gives a wrong answer. If two independent sensors disagree, one is failing, so we
don't average.}

\qa{Why 8 BPM?}
{Slightly above the AAMI tolerance ($\pm5$\,BPM) to allow for each channel's own error. It is a tunable parameter, to be
re-tuned on development data.}


% ---------------------------------------------------------------
\qgroup{H. Colour indices (Diagnostics Studio)}

\qa{Can a camera detect jaundice or anaemia?}
{It can give \textbf{screening hints}. Research (e.g.\ BiliScreen) used colour-calibration cards and controlled light. Ours are
\textbf{uncalibrated and exploratory}; lighting and skin tone change them a lot.}

\qa{Your screenshot says jaundice ALERT. Do you have jaundice?}
{\textbf{Trap.} No: warm room light lowers blue, and the eye band includes skin and eyebrows, which raises the yellowness
index. It shows exactly why calibration is needed.}

\qa{Why log in the erythema index?}
{By the Beer--Lambert law, absorbance is logarithmic in reflectance, so $\log R - \log G$ tracks haemoglobin absorption roughly
linearly.}

\qa{What is the Laplacian variance?}
{A sharpness measure: the variance of the 2nd-derivative image. Blurred frames have weak edges and a low value; we pick the
sharpest frame for colour analysis.}

\qa{How would you make colour indices reliable?}
{A colour reference card in the frame, fixed lighting, white-balance correction, and validation against lab tests (bilirubin,
haemoglobin) with ethics approval.}


% ---------------------------------------------------------------
\qgroup{I. Validation, accuracy and limitations}

\qa{How accurate is your heart rate?}
{\textbf{Not yet validated} against a reference device; we will not claim a number we haven't measured. Validated so far: the
algorithms on synthetic signals (31 automated tests), and the streaming link (8\,fps, 0 loss). Planned: ECG chest strap as
reference.}

\qa{How will you validate it?}
{Simultaneous ECG chest strap (HR, R--R intervals), finger oximeter, metronome breathing; conditions at rest, talking, small
motion, low light, different distances; 30\,s windows; \textbf{MAE, RMSE, Pearson $r$}, \textbf{Bland--Altman} bias and 95\,\% limits
of agreement, and the \textbf{AAMI EC13} criterion (error $\leq 5$\,BPM or 10\,\%).}

\qa{What is a Bland--Altman plot?}
{Difference (device $-$ reference) against their mean. The \textbf{bias} is the mean difference; the \textbf{limits of agreement}
are bias $\pm1.96$\,SD. It shows agreement; correlation only shows that two values move together. With several windows per
person, the repeated-measures version is used.}

\qa{Why not just use correlation?}
{Two methods can correlate perfectly yet differ by 10\,BPM. Agreement needs bias and limits (Bland \& Altman, 1986).}

\qa{Why split development and test participants?}
{Thresholds (SNR gate, $\Delta$, RR gate) are tuned on development participants only, so test results are not biased by
tuning.}

\qa{Main limitations?}
{\textbf{Motion} and \textbf{lighting changes}; low resolution, 8\,fps and JPEG compression; \textbf{skin tone} (more melanin means a
weaker pulse signal, a known bias in rPPG); beard and hair covering ROIs; weak heart sounds through air; no clinical
validation yet.}

\qa{Does skin tone matter?}
{Yes. Melanin absorbs light and reduces the pulse amplitude, so SNR is lower for darker skin (Nowara et al., 2020). The SNR gate
avoids wrong numbers but may show ``--'' more often. Diverse participants are needed.}

\qa{What about fluorescent light flicker?}
{Mains lighting flickers at 100\,Hz. With 8\,fps sampling it can alias into the low-frequency band as a fake peak. Locking the
exposure and using daylight or DC LED light reduce it.}

\qa{What if I show it a photo of a face?}
{The face detector will find it, but a photo has no pulse: the spectrum is noise, the SNR stays below the gate, and HR shows
``--''. (A photo under flickering light could still pass occasionally; liveness detection is future work.)}

\qa{Why did your earlier screenshot show 97 BPM and another 57 BPM?}
{Different moments; and in that first build there was no quality gate, so noisy windows were still displayed. The current
build only shows estimates that pass the SNR gate.}

\qa{Your HRV was exactly 45.0 and stress exactly 50.0. Why?}
{\textbf{Trap; be honest.} In the first build those were placeholder formulas driven by the quality score; with quality at its
floor (50) they gave 45.0 and 50.0. I found this, removed them, and replaced HRV with real RMSSD from heart sounds.}

\qa{What is the latency?}
{About 10\,s for the first estimate (buffer), then an update every 2\,s on the dashboard and every 1.5\,s to the OLED.}


% ---------------------------------------------------------------
\qgroup{J. Software}

\qa{Software stack?}
{Firmware: Arduino core on ESP-IDF/FreeRTOS, built with PlatformIO. Host: Python, Flask, OpenCV, NumPy/SciPy. Frontend:
HTML/CSS/JavaScript with Plotly charts, polling the Flask JSON API.}

\qa{Why process on the laptop instead of on the ESP32?}
{Face detection, CHROM and spectral analysis in Python are easy to develop, debug and change. The ESP32 could run a simpler
version (e.g.\ fixed ROI plus FFT): \textbf{edge processing} is future work and would also remove the video stream.}

\qa{How do you record data for analysis?}
{The ``Experiments'' tab starts a session: \texttt{frames.csv} (timestamps, ROI colour, gate), \texttt{audio.wav},
\texttt{estimates.csv}, \texttt{meta.json}. \texttt{analysis/evaluate.py} replays them offline with the \textbf{same} DSP code.}

\qa{How did you test the software?}
{31 automated \texttt{pytest} tests with synthetic signals of known HR/RR (e.g.\ 52/78/128\,BPM within 2\,BPM), protocol parsing,
gap bridging, noise rejection, and a node simulator that imitates the ESP32 stream.}

\qa{Privacy and security?}
{Everything stays on the local network; nothing goes to the cloud; video is saved only during an explicit recording. Limitation:
the TCP stream is not encrypted (WPA2 only protects the Wi-Fi link).}


% ---------------------------------------------------------------
\qgroup{K. Future work}

\qa{What would you improve first?}
{(1)~A validation study with an ECG reference; (2)~compute ROI colour means \textbf{on the ESP32 from raw pixels} before JPEG
(avoids compression loss); (3)~a stethoscope-coupled or contact mic; (4)~a landmark face tracker; (5)~an IR camera/LED for SpO$_2$;
(6)~edge processing and a battery.}

\qa{Could this become a product?}
{After clinical validation, IEC 60601 electrical-safety testing and regulatory approval (CDSCO in India) as a medical device. Today
it is a research prototype.}

\end{multicols}
